\begin{resumen}

El problema de la satisfacibilidad booleana (SAT) se estudia en este
trabajo desde la teoría de lenguajes formales. En la línea de
investigación en la que se inscribe, la satisfacibilidad de una fórmula
se reduce al vacío de la intersección de un lenguaje regular de
interpretaciones con una gramática que impone la consistencia entre las
ocurrencias de cada variable booleana.
Con gramáticas menos expresivas esa reducción se limita a fórmulas con
cierto orden de las ocurrencias: una gramática libre del contexto no
admite cruces, y en una de concatenación de rango simple (sRCG) la aridad
crece con la fórmula. En un intento de resolver más fórmulas en tiempo
polinomial se ataca el problema con gramáticas de concatenación de rango
(RCG): con su no linealidad se modela la consistencia de más fórmulas y
se mantiene mínima la aridad, como se muestra con una RCG que reconoce la
consistencia de cualquier fórmula de SAT con aridad uno; pero por esa
misma no linealidad se pierde la garantía directa de que el vacío sea
polinomial.
Se da un algoritmo de fuerza bruta que decide ese vacío: se enumeran las
cadenas del autómata y se reconocen con la RCG; su costo es exponencial,
así que
no es una vía factible, pero muestra por qué conviene estudiar el
problema sobre un único objeto. Se construye ese objeto, la gramática
producto, y se caracterizan su tamaño y el costo de construirlo; sobre él
se puede abordar el vacío: hallar un método polinomial para todas las
fórmulas, lo que implicaría $P = NP$; probar que no existe; o hallarlo al
restringir a una subclase. Aunque su
tamaño no es polinomial, queda como objeto de estudio: eliminar sus
cláusulas no productivas y analizar el costo de decidir el vacío se
proponen como trabajo futuro.

\vspace{1em}
\noindent\textbf{Palabras clave:} satisfacibilidad booleana, gramáticas
de concatenación de rango, vacío de la intersección, lenguaje regular
finito, no linealidad.

\end{resumen}

\begin{abstract}

The Boolean satisfiability problem (SAT) is studied in this work from
formal language theory. Within the line of research to which it belongs,
the satisfiability of a formula reduces to the emptiness of the
intersection of a regular language of interpretations with a grammar that
enforces consistency among the occurrences of each Boolean variable. With
less expressive grammars this reduction is limited to formulas with a
certain order of occurrences: a context-free grammar admits no crossings,
and with a simple range concatenation grammar (sRCG) the arity grows with
the formula. In an attempt to solve more formulas in polynomial time, the
problem is approached with range concatenation grammars (RCG): with their
non-linearity the consistency of more formulas is modeled while the arity
is kept minimal, as shown by an RCG that recognizes the consistency of
any SAT formula with arity one; but with that same non-linearity the
direct guarantee that the emptiness is polynomial is lost. A brute-force
algorithm that decides this emptiness is given: the strings of the
automaton are enumerated and recognized with the RCG; its cost is
exponential, so enumeration is not a feasible route, but it shows why the
problem is better studied on a single object. This object, the product
grammar, is constructed, and its size and the cost of constructing it are
characterized; on it the emptiness can be addressed: finding a polynomial
method for all formulas, which would imply $P = NP$; proving that none
exists; or finding one only by restricting the modeling to a subclass.
Although its size is not polynomial, it remains an object of study:
removing its non-productive clauses and analyzing the cost of deciding
the emptiness are proposed as future work.

\vspace{1em}
\noindent\textbf{Keywords:} boolean satisfiability, range concatenation
grammars, emptiness of the intersection, finite regular language,
non-linearity.

\end{abstract}
